\documentclass[11pt,letterpaper]{article}
\usepackage[T1]{fontenc}
\usepackage{lmodern,amsmath,amssymb,geometry,xcolor,hyperref}
\geometry{left=.94in,right=.94in,top=.77in,bottom=.83in}
\definecolor{ink}{HTML}{19343A}
\definecolor{teal}{HTML}{126267}
\hypersetup{colorlinks=true,urlcolor=teal,linkcolor=teal,pdftitle={When does multinomial averaging preserve monotonicity?},pdfauthor={LDJAI with OpenAI Codex}}
\setlength{\parindent}{0pt}
\setlength{\parskip}{5pt plus .5pt}
\setlength{\abovedisplayskip}{6pt plus 1pt minus 1pt}
\setlength{\belowdisplayskip}{6pt plus 1pt minus 1pt}
\newcommand{\E}{\mathbb E}
\newcommand{\R}{\mathbb R}
\newcommand{\Z}{\mathbb Z}
\newcommand{\head}[1]{\par\vspace{5pt}{\large\bfseries\color{ink}#1}\par\vspace{1pt}}
\newcommand{\subhead}[1]{{\bfseries\color{ink}#1}\quad}
\newcommand{\gapF}{\mathcal D_n}
\begin{document}
{\small\sffamily\bfseries\color{teal} PROBABILITY}\hfill{\small November 2026}

{\LARGE\bfseries\color{ink} When does multinomial averaging\\[2pt] preserve monotonicity?\par}
{\small LDJAI with OpenAI Codex}

Let $p$ lie in the probability simplex $\Delta_{k-1}$, and let $X_p$ count the outcomes of $n$ independent categorical draws with probabilities $p$. For a field $f\colon\Z^k\to\R^k$, put $G_n(p)=\E[f(X_p)]$. In 2010, daizhuo asked whether monotonicity of $f$ in the sense of
\[
 (x-y)\cdot(f(x)-f(y))\ge0 \qquad(x,y\in\Z^k)
\]
implies the same inequality for $G_n$ on the simplex. Write
$\gapF(p,q)=(p-q)\cdot(G_n(p)-G_n(q))$.

\subhead{Theorem.} For every $n\ge1$, multinomial averaging preserves monotonicity of every field on $\Z^k$ if and only if $k\le2$. For $k\ge3$, a single field and a fixed pair of probability vectors violate the inequality for every $n\ge1$.

\head{Three categories suffice to reverse the sign}
For $k=3$, take $a=(-1,0,1)$ and $f(x)=a\,(a\cdot x)_+$, where $t_+=\max(t,0)$. For any $x,y$, the monotonicity inequality reduces to
\[
 (a\cdot x-a\cdot y)\big((a\cdot x)_+-(a\cdot y)_+\big)\ge0,
\]
since the positive part is increasing. Compare $p=(0,1,0)$ with $q=(3/4,0,1/4)$. Under $p$ the score $a\cdot X_p$ is zero. Under $q$ it equals $n$ whenever every draw selects category three, an event with probability $4^{-n}$. Consequently
\[
 G_n(p)=0,\qquad G_n(q)=a\,\E[(a\cdot X_q)_+],\qquad
 \gapF(p,q)=-\tfrac12\E[(a\cdot X_q)_+]\le-\frac{n}{2\cdot4^n}<0.
\]
The mean score drops from $0$ to $-n/2$, yet the expected positive part increases. For $k>3$, append zeros to $a,p,q$. The same monotone field formula and the same inequality hold.

\head{Two categories reduce to one ordered count}
At $k=2$, write $p_t=(t,1-t)$ and $z_j=(j,n-j)$. The scalar contrast $h(j)=f_1(z_j)-f_2(z_j)$ is nondecreasing in $j$, by applying monotonicity to $z_{j+1},z_j$ whose difference is $(1,-1)$. On one probability space, let
$J_t=\sum_{i=1}^n\mathbf{1}_{\{U_i\le t\}}$ for independent $U_i$ uniform on $[0,1]$. Then $z_{J_t}$ has the required multinomial law, and $J_t$ increases with $t$ for every realization. It follows that
\[
 \gapF(p_t,p_s)=(t-s)\big(\E h(J_t)-\E h(J_s)\big)\ge0.
\]
For $k=1$ the simplex has one point, completing the proof.

\newpage
\head{A smooth, strongly monotone counterexample}
The sign reversal persists for the real-analytic gradient of a strongly convex potential. Put $a=(0,1,2)$, $p=(3/5,0,2/5)$, $q=(0,1,0)$, and define on $\R^3$
\[
 F(x)=a\,2^{a\cdot x}+\frac{x}{76}
 =\nabla\!\left(\frac{2^{a\cdot x}}{\log 2}+\frac{\lVert x\rVert_2^2}{152}\right).
\]
The Jacobian $DF(x)=(\log 2)2^{a\cdot x}aa^{\mathsf T}+I/76$ has all eigenvalues at least $1/76$ and all entries nonnegative. Thus $F$ is $1/76$-strongly monotone, is nondecreasing in every coordinate, and takes componentwise nonnegative values on the nonnegative orthant. Restricting $F$ to $\Z^3$ gives a field of the type in the conjecture.

Let $z(r)=\sum_i r_i2^{a_i}$. The independence of the draws gives $\E[2^{a\cdot X_r}]=z(r)^n$, while $\E[X_r]=nr$. For $G_n(r)=\E[F(X_r)]$ this yields
\[
 G_n(r)=a\,z(r)^n+\frac{nr}{76}.
\]
Here $z(p)=11/5$, $z(q)=2$, $(p-q)\cdot a=-1/5$ and $\lVert p-q\rVert_2^2=38/25$. The difference-of-powers identity gives $(11/5)^n-2^n\ge n/5$, so
\[
 \gapF(p,q)=-\frac15\big((11/5)^n-2^n\big)+\frac{n}{50}
 \ \le\ -\frac{n}{50}\qquad(n\ge1).
\]
The violation survives with strictly positive probabilities, using the same pair for all $n$. Fix $0<\varepsilon<1$ and any interior probability vector $u$. Replace $p,q$ by $p'=(1-\varepsilon)p+\varepsilon u$ and $q'=(1-\varepsilon)q+\varepsilon u$. Then $z(p')-z(q')=(1-\varepsilon)/5$ with both bases at least one. Hence $z(p')^n-z(q')^n\ge n(1-\varepsilon)/5$. Both the resulting negative upper bound on the exponential contribution and the positive quadratic contribution scale by $(1-\varepsilon)^2$, giving
\[
 \gapF(p',q')\le-\frac{n(1-\varepsilon)^2}{50}.
\]

In a model that evaluates a monotone cost field on sampled counts, its expected cost can therefore violate monotonicity as a function of the sampling probabilities.

\head{Two cases suggested in the original question}
For affine $f(x)=Ax+b$, one has $G_n(p)=nAp+b$. Lattice monotonicity gives $v^{\mathsf T}Av\ge0$ first for integral $v$, then for rational $v$ by scaling and for real $v$ by continuity. Hence $\gapF(p,q)\ge0$. If $f_i(x)=h_i(x_i)$ with each $h_i$ nondecreasing, $X_{p,i}$ is binomial with parameter $p_i$. The uniform coupling on page one shows that $\E[h_i(X_{p,i})]$ increases with $p_i$. Every term in $\gapF$ is then nonnegative.

\head{Source and verification}
\small The conjecture is from daizhuo, \emph{Expectation, multinomial distribution, and monotonicity (A conjecture)}, MathOverflow, December 1, 2010, \href{https://mathoverflow.net/questions/47888/expectation-multinomial-distribution-and-monotonicity-a-conjecture}{question 47888}. The \href{https://lefttosolve.com/p/mo-47888}{Left to Solve catalog entry} is \texttt{mo-47888}. Prepared with OpenAI Codex at LDJAI's request. The companion Lean 4 files verify the counterexample for all $n\ge1$ and the full category classification. The smooth example is proved above.
\newpage
\normalsize
{\small\sffamily\bfseries\color{teal} APPLICATION TO ROUTE CHOICE}
\head{Which linear load measurements preserve monotonicity?}
Let $A$ be a real $m\times k$ matrix. For route or activity counts $x$, the vector $Ax$ records $m$ resource loads. Assign each load a nondecreasing scalar cost $c_e\colon\R\to\R$ and put $f(x)=A^{\mathsf T}c(Ax)$, with $c(v)_e=c_e(v_e)$. This field is monotone because its defining inner product is a sum of one-variable monotonicity inequalities.

\subhead{Theorem.} For any fixed $n\ge1$, the expected field $G_n(p)=\E[f(X_p)]$ is monotone for every choice of the $c_e$ if and only if each row of $A$ has at most two distinct entries.

\subhead{Sufficiency.} For a row $u$ with entries in $\{a,b\}$, $a<b$, let $t_p=\sum_{i:u_i=b}p_i$. Each draw contributes either $a$ or $b$, so
\[
 u\cdot X_p\ \overset{d}{=}\ na+(b-a)J_{t_p},\qquad J_t\sim\mathrm{Binomial}(n,t).
\]
The expected cost $H(t)=\E[c_e(na+(b-a)J_t)]$ is nondecreasing by the coupling on page one. Since $u\cdot(p-q)=(b-a)(t_p-t_q)$, the contribution of this row to $\gapF(p,q)$ equals $(b-a)(t_p-t_q)(H(t_p)-H(t_q))\ge0$. A constant row contributes zero. Summing over rows proves sufficiency, even for correlated resource loads.

\subhead{Necessity.} Suppose one row $u$ takes values $l<m_0<h$. Set its cost to $\exp$ and all others to zero. Let $q$ select the middle value $m_0$ with certainty. Choose an endpoint mixture $p_0$ with mean $m_0$. Strict convexity gives $\E_{p_0}[e^u]>e^{m_0}$ for one draw. Move a little mass from $h$ to $l$ to obtain $p$. Its mean load per draw falls below $m_0$ while its mean exponential stays above $e^{m_0}$. With $v_i=e^{u_i}$, independence yields
\[
 \gapF(p,q)=\underbrace{u\cdot(p-q)}_{<0}
      \underbrace{\big((v\cdot p)^n-(v\cdot q)^n\big)}_{>0}<0.
\]
Mixing $p,q$ with the same interior distribution in the same proportion preserves both strict orders. Thus one strictly positive pair witnesses failure for every $n\ge1$.

\head{Expected route costs on a network}
Let $A_{ei}$ indicate whether route $i$ traverses link $e$. A link with use probability $(Ap)_e$ then has load $\mathrm{Binomial}(n,(Ap)_e)$. Write $h_e(t)=\E[c_e(J_t)]$. For travelers who choose routes independently, the expected vector of additive route costs is $A^{\mathsf T}h(Ap)$, and
\[
 (p-q)\cdot(G_n(p)-G_n(q))=
 \sum_e\big((Ap)_e-(Aq)_e\big)\big(h_e((Ap)_e)-h_e((Aq)_e)\big)\ge0 .
\]
For example, if routes 1 and 2 share a link of cost $j^2$ at load $j\ge0$ and route 3 bypasses it, the shared link adds $(x_1+x_2)^2$ to the costs of both routes. Its expected contribution to each is $nt(1-t)+n^2t^2$ for $t=p_1+p_2$. The route costs can therefore be nonlinear and depend on counts of other routes while retaining expected monotonicity.

{\small\textbf{Traffic-model motivation.} Daizhuo Chen, \emph{Modeling travel time uncertainty in traffic networks}, S.M. thesis, Massachusetts Institute of Technology, 2010.
\href{https://dspace.mit.edu/handle/1721.1/61889}{MIT thesis record}. The real-matrix criterion above is proved directly on this page.}
\end{document}
